% 原文 pp.195–201（印刷 pp.1313–1319）
\section{Lecture 14：二维 $\mathcal{N}=2$ 超对称理论（二）：手征环与扭曲理论}

\begin{keytakeaways}
本讲在二维 $\mathcal{N}=2$ 超对称理论中取定一个超荷 $Q$，把手征环（chiral ring）定义为局部算符的 $Q$ 上同调，并证明 $Q$-闭算符的「后代」在闭链上的积分给出与度规无关、因而具有拓扑意义的关联函数。
核心数学工具是 $R$ 对称性的规范化与「扭曲」：把旋转生成元改成 $K'=K\pm\frac12J_L\pm\frac12J_R$，使 $Q$ 在任意超 Riemann 曲面上都是整体超对称；可规范化的组合只有 $J_L\pm J_R$，其判据是反常系数二次型在这两个零向量上为零（$k=0$）。
物理图像是：扭曲之后理论成为拓扑场论，路径积分只依赖 $Q$ 的不动点集——纯 $\mathcal{N}=2$ 规范理论给出平坦联络模空间上的示性类积分，Kähler $\sigma$ 模型给出量子上同调中的积分。
主要结论是：$Q$ 上同调类不依赖插入点的位置，故手征环是 $\mathbb{Z}/2\mathbb{Z}$ 分次交换环；在 Hamilton 框架下 $H=\{Q,\widetilde Q\}$、$H^*(Q)=\ker Q/\operatorname{Im}Q$，当 $H$ 有离散谱时它等同于零能态空间。
\end{keytakeaways}

% ===== 原文 p.195 =====
% 原文页眉：LECTURE 14；标题：$N=2$ SUSY Theories in Dimension Two, Part II: Chiral Rings and Twisted Theories

本讲再次讨论二维的 $\mathcal{N}=2$ 超对称理论。
取其中一个超对称 $Q$，我们将对定义在平坦闵可夫斯基空间或欧氏空间中的理论，
定义局部算符的手征环（chiral ring）与态的 $Q$ 上同调（$Q$-cohomology）。
接着我们讨论如何扭曲（twist）理论，使超对称 $Q$ 至少在一大类 $2|4$ 维超流形（supermanifold）上具有整体意义。
一旦 $Q$ 在 $\Sigma$ 上的某个理论中被整体化（globalize），我们便有了该理论的「手征局部算符环」这一概念。
我们可以定义所谓 $Q$ 不变算符的后代（descendants），
并证明：这些后代在流形中各闭链上的积分给出具有拓扑意义的关联函数。
为计算这些函数，人们对 $Q$ 不变的场作积分。

随后我们详细考察两个带整体超对称 $Q$ 的扭曲理论的例子。
第一个是二维的、规范群为单 Lie 群 $G$ 的纯 $\mathcal{N}=2$ 规范理论。
在这里我们看到，手征函数（chiral functions）之环被等同于 Lie 代数上的不变多项式（invariant polynomial）。
经典地，$Q$ 不变态的空间就是通常的平坦联络（flat connection）空间
（在联络可约（reducible）的情形，还要加上伴随丛（adjoint bundle）的一个额外平行截面）。
$Q$ 不变算符的关联函数被等同于：由 Lie 代数中的不变多项式导出的那些通常类，
在平坦联络模空间上的积分。（可约联络会带来额外的复杂性。）

第二个例子是映入 Kähler 流形的 $\mathcal{N}=2$ 超对称二维 $\sigma$ 模型。
结果表明，存在一种扭曲方式，它给出的手征环正是该 Kähler 流形的拓扑上同调（topological cohomology）。
$Q$ 不变态的空间是从 Riemann 曲面到 $X$ 的全纯映射（holomorphic map）空间。
在这个扭曲理论中，计算拓扑关联函数的路径积分正是量子上同调（quantum cohomology）中的积分。
当 Kähler 流形是 Calabi--Yau 流形时，还有第二种不等价的扭曲方式，它给出一个不同的理论。
在该理论中，手征环是 $H^*(X,\wedge^*T^{1,0}(X))$，
而 $Q$ 不变的场是从 Riemann 曲面到 $X$ 的常映射。
于是，计算拓扑关联函数的路径积分就是这些对象在 $X$ 上的积分——
这些积分与 $X$ 上 Hodge 结构的形变（variation of Hodge structure）有关。

\noindent\textit{（本讲笔记由 John Morgan 整理。原文 p.1313 页脚）}

% [Tutor Note: 这里的 $H^*(X,\wedge^*T^{1,0}(X))$ 用的是全纯**切**丛 $T^{1,0}$ 的外幂，而不是通常写 Hodge 分解时用的全纯余切丛 $\Omega^p=\wedge^pT^{*(1,0)}$，抄写时极易在这一步出错。对 Calabi--Yau 流形，$K_X$ 平凡给出 $\wedge^pT^{1,0}\cong\Omega^{\,n-p}$（$n=\dim_{\mathbb C}X$），故二者作为分次向量空间的总和相同，只是标号 $p\leftrightarrow n-p$ 互换——这正是作者说这些积分「与 Hodge 结构的形变有关」的原因。另可注意：本页先后描述的两种扭曲即今称的 A 模型（手征环 $=H^*(X)$，态 $=$ 全纯映射）与 B 模型（手征环 $=\oplus_{p,q}H^q(X,\wedge^pT^{1,0})$，态 $=$ 常映射），二者的差别仅是 $K'$ 中 $J_R$ 前的符号。]

% [Tutor Note: 原文此处写 "the ring of chiral functions"，按上下文（同页前面说的是 chiral ring of local operators）应读作「手征算符之环」；"functions" 一词沿用了作者在讨论超势时对「超场的函数」的习惯说法，并非指普通函数。译文按「手征函数之环」直译以免改动原文，但读者须知它指的是算符的环。]

% ===== 原文 p.196 =====
\subsection{$R$ 对称性再论}
% 原文 §14.1：R-symmetry revisited

我们所考察的大多数 $\mathcal{N}=2$ 超对称例子都将（至少在经典层面）同时具有左行与右行的 $R$ 对称性——
即上一讲中记作 $J_L$ 与 $J_R$ 的那些 $R$ 对称性。
只用其中一个 $R$ 对称性也可以完成今天的讨论，但至少必须有一个。
不过，我们仍将假定 $R$ 对称性群是 $U(1)\times U(1)$，从而 $J_R$ 与 $J_L$ 兼备。
理论定义所在的超流形是 $\mathbb{R}^{2|4}$。
我们最终还将允许某一类带欧氏度量的分裂（split）$2|4$ 维超 Riemann 曲面
（以及带平坦闵可夫斯基度量的 $2|4$ 维超流形）。
奇方向的空间是 $J_L$ 的二维表示的一份拷贝与 $J_R$ 的二维表示的一份拷贝。
我们预期 $J_L$ 与 $J_R$ 之间存在由世界面定向反转给出的对称性。

我们可以对 Lagrange 量的一个对称性尝试做的一种操作，是把它规范化（gauge）。
若待规范化的对称群是 $H$，则新的 Lagrange 量将含有规范群为 $H$ 乘以原理论规范群（若有）的规范场。
直观上，规范化过程就是把取值于一个常空间的场，换成取值于一个 $H$-丛的场，并添加一个新场——$H$-联络。
当作用量或 Lagrange 量是某个局部 Lagrange 密度的积分时，规范化最为直接。
此时在经典层面总可以构造出新的、已规范化的 Lagrange 量：把 Lagrange 量中的普通导数换成协变导数。
往往存在一种自然的做法来完成这一过程；但无论如何，对对称性作规范化的不同选择
所得到的 Lagrange 量之间只差规范不变项。
然而还有一个问题：规范化后的理论作为量子理论是否有意义。
若原理论具有这一性质，那么（在我们所处维数足够低、不会因规范化而遭遇新的紫外困难的前提下）
规范化后的理论也具有该性质的充要条件是：某些反常（anomaly）必须相消。

在我们的情形中，永远不可能把整个 $U(1)\times U(1)$ 的 $R$ 对称性群规范化，
而得到一个作为量子理论有意义的结果。
事实上，能以这种方式规范化的 $J_L$ 与 $J_R$ 的组合只有 $J_L\pm J_R$。
它们被称为矢量流与轴矢流：
\[
\begin{aligned}
J_V &= J_L+J_R,\\
J_A &= J_L-J_R.
\end{aligned}
\]
若其中之一被规范化，另一个就会出现反常，不再是量子理论的对称性。
为说明这一点，我们需要关于反常的一些一般性观察。
说一个理论具有由流 $J_L$ 生成的对称性，意思是关联函数满足恒等式
\[
d_x\bigl\langle J_L(x)\mathcal{O}(y_1)\dots\bigr\rangle=0 \qquad (\ast)
\]
只要 $x$ 与各 $y_i$ 不同。这里 $d_x$ 是对 $x$ 变量的外微分。
若把关联函数（作为分布）延拓穿过对角线，这个恒等式会多出支撑在对角线上的项。
一个非常重要的特例是流本身的两点函数（或者，在 $2k$ 维时空中是 $(k+1)$ 点函数）。
对此，一般结构是
\[
d_x\bigl\langle J_L(x)J_L(y)\bigr\rangle=\frac{k}{2\pi}d_y\delta(x-y) \qquad (\ast\ast)
\]

% [Tutor Note: $(\ast\ast)$ 中的系数 $\frac{k}{2\pi}$ 依赖流的归一化，切勿与「级」（level）混为一谈：若把 $J_L$ 重新标度 $J_L\to cJ_L$，则 $k\to c^2k$。真正有物理意义、且在下页被用到的是 $k$ 的整性（p.197 所谓「满足一个整性条件」）；在二维共形场论中它就是 $U(1)$ 流的量子化级，故「$k=0$ 才可规范化」等价于「级为零才无反常」。另请注意作者用「$2k$ 维时空中的 $(k+1)$ 点函数」一句概括了一般维数的反常：$dJ$ 的对角项在四维表现为三点函数（三角形图），在二维表现为两点函数，后者正是下一段 $dJ_L=\frac{k}{2\pi}F$ 的来源。]

% [Tutor Note: 原文作 "The supermanifold on which the theory is defined in $\mathbf{R}^{2|4}$"，其中的介词 "in" 显然是 "is" 之误（试与下一句 "We shall also eventually allow…" 对照）。译文按「理论定义所在的超流形是 $\mathbb{R}^{2|4}$」处理；此处若照字面读成「定义在 $\mathbb{R}^{2|4}$ 中」虽不影响物理，却会漏掉作者正在引入「超流形」这一几何框架的用意。]

% ===== 原文 p.197 =====
其中 $k$ 是某个常数，它可以为零，也可以不为零。
生成 $J_L$ 的整体对称性可被规范化的条件，恰恰是 $k=0$。
为看清这一点，让我们试着通过向理论中添加一个 $U(1)$ 规范场 $A$ 来规范化该对称性。
Lagrange 量获得一个额外项 $\int A\wedge J_L$。
（若 $J_L$ 本身含有带荷场的导数，还可能有更多项，即 $A$ 的高阶项。）
现在，我们要判断带这个耦合的理论是否真的具有由 $J_L$ 生成的对称性，
即恒等式 $(\ast)$ 对这个理论是否成立。
为此我们必须在原理论中考察
\[
d_x\left\langle J_L(x)\mathcal{O}(y_1)\dots\exp\left(\int A\wedge J_L\right)\right\rangle.
\]
引入指数是为了用原理论来表达新理论的关联函数。
若把它按 $A$ 的幂展开，则线性于 $A$ 的项在利用 $(\ast\ast)$ 并分部积分之后给出
\[
\bigl\langle \mathcal{O}(y_1)\cdots kF(x)\bigr\rangle,
\]
其中 $F=dA$ 是曲率。这可以写成一个算符关系
\[
dJ_L=\frac{k}{2\pi}F,
\]
它断言：规范化之后，$J_L$ 不再生成对称性。

更一般地，若有若干流 $J^1,\dots,J^s$ 生成 $s$ 维的 Lie 群 $G$，其 Lie 代数为 $\mathcal{G}$，
则两点函数满足
\[
d_x\bigl\langle J^a(x)J^b(y)\bigr\rangle=\frac{k^{ab}}{2\pi}d_y\delta(x-y),
\]
其中 $k^{ab}$ 是 $\mathcal{G}$ 上某个不变二次型。
二次型 $k$ 没有正定性（例如左行与右行无质量费米子的贡献符号相反），但确实满足一个整性条件。
若我们规范化 $G$ 的一个子群 $H$，其 Lie 代数为 $\mathcal{H}$，
则按与上面类似的推理，量子理论只有在二次型 $k$ 限制到 $\mathcal{H}$ 上为零时才有意义。
若确实如此，则规范化后理论的「整体对称性」由 $\mathcal{H}^{\perp}/\mathcal{H}$ 中的流生成，
也就是 $\mathcal{G}$ 中与 $\mathcal{H}$ 正交的流，模去 $\mathcal{H}$ 中的那些流。

在我们将要考察的 $\mathcal{N}=2$ 理论中有两个流 $J_L,J_R$，
并且存在一个交换 $J_L$ 与 $J_R$ 的「宇称」对称性（即二维时空的定向反转）。
$k$ 在定向反转下为奇，故在这种情形下 $k$ 的结构是
$k^{LL}=-k^{RR}\neq0$、$k^{LR}=0$。
这个二次型有两个零向量 $J_L\pm J_R$，
因此这两个 $R$ 对称性的线性组合中任何一个都可以被规范化。
若其中一个被规范化，第二个就变得「反常」——它不再守恒——
因为这两个零向量当然并不彼此正交。
「幸运的是」，这些选项与我们在下面因其他理由所需要的选择一致。

我们真正感兴趣的并不只是规范化，而是用 $R$ 对称性来扭曲理论。
这一过程要求先规范化 $R$ 对称性，然后用一个从度规所决定的正交群到对称群的同态来扭曲它。
结果表明，进行这一过程将有两种等价的方式。
但在着手扭曲之前，我们必须先理解局部的情形。那是下一节的主题。

% [Tutor Note: 「两个零向量并不彼此正交」可以一行算出：设 $k^{LL}=k_0$、$k^{RR}=-k_0$、$k^{LR}=0$，则 $k(J_L+J_R,J_L-J_R)=k_0+k_0=2k_0\neq0$，而 $k(J_L\pm J_R,J_L\pm J_R)=k_0-k_0=0$，故 $J_L\pm J_R$ 都是零向量、却互不正交。正是这个非零内积使得「规范化其中一个 $\Rightarrow$ 另一个反常」：规范场 $A_\pm$ 的混合 Chern--Simons 型项（在二维就是 $A_+A_-$ 型的质量项）会破坏另一条流的守恒。作者说「幸运的是这些选项与下面所需一致」，指的是 §14.3 中要规范化的恰是 $K'=K\pm\frac12J_L\mp\frac12J_R$ 所对应的组合，即 $J_L$ 与 $J_R$ 的同号/反号线性组合。]

% [Tutor Note: $\mathcal{H}^{\perp}/\mathcal{H}$ 这一写法容易读错：这里「正交」是对不变二次型 $k$ 而言的，而非对 Killing 形式；当 $k$ 在 $\mathcal{H}$ 上退化时（本讲正是如此），$\mathcal{H}^\perp$ 严格大于 $\mathcal{H}$，商空间才给出真正的「剩余整体对称性」。另注意 $\mathcal{H}^{\perp}$ 必须理解为 $\{X\in\mathcal{G}\,|\,k(X,Y)=0\ \forall Y\in\mathcal{H}\}$，而不是 $\mathcal{G}$ 中与 $\mathcal{H}$ 对易的子代数——两种「正交」在这里给出不同的答案。]

% ===== 原文 p.198 =====
\subsection{算符的 $Q$ 上同调}
% 原文 §14.2：Q-cohomology of operators

我们考虑带标准欧氏或闵可夫斯基度量的 $\mathbb{R}^{2|4}$。
设我们的二维 $\mathcal{N}=2$ 超对称理论具有如下超对称性：
\begin{center}
\begin{tabular}{c|c}
\hline
$U(1)_L$ 的特征标 & $U(1)_R$ 的特征标\\
\hline
$Q_-$ & $Q_+$\\
$\overline{Q}_-$ & $\overline{Q}_+$\\
\hline
\end{tabular}
\end{center}
$Q_-$ 与 $\overline{Q}_-$ 在 $U(1)_L$ 的标准表示及其共轭表示中变换，并与 $U(1)_R$ 对易；
而 $Q_+$、$\overline{Q}_+$ 与 $U(1)_R$ 及 $U(1)_L$ 的关系类似。
这些超对称满足关系
\[
Q_+^2=\overline{Q}_+^2=0,\qquad \{Q_+,\overline{Q}_+\}=P_+,
\]
以及把 $+$ 换成 $-$ 的对称关系。此外，
\[
\{Q_+,Q_-\}=\{\overline{Q}_+,Q_-\}=\{Q_+,\overline{Q}_-\}=\{\overline{Q}_+,\overline{Q}_-\}=0.
\]
我们将同时考虑四种密切相关的可能性。令 $Q$ 为下列之一：
\[
Q_++Q_-,\qquad Q_++\overline{Q}_-,\qquad \overline{Q}_++Q_-,\qquad \overline{Q}_++\overline{Q}_-.
\]
（前两种可能性给出不等价的扭曲理论；后两种是前两种的厄米共轭，会给出等价的理论。）
为固定记号，我们这样选取生成元，使得 $Q=Q_++Q_-$。
令 $\mu=\pm$，我们看到 $P_\mu$ 可以写成一个对易子
\begin{equation}
  \tag{14.1}
  \label{eq:14-1}
  P_\mu=\{Q,T_\mu\},
\end{equation}
其中 $T_\mu$ 是 $Q_\pm,\overline{Q}_\pm$ 的一个线性组合。
我们感兴趣的是 $Q$ 作为微分（通过取超括号（superbracket）作用）在局部算符空间上
（更一般地，作用在这类算符的乘积上）的上同调。
于是，一个局部算符 $\phi(x)$ 若满足 $Q$-闭，即 $\{Q,\phi(x)\}=0$，就定义一个 $Q$ 上同调类。
用式 14.1 作一个简单计算即可看出，对 $Q$-闭算符 $\phi(x)$ 我们有
\[
d_\mu\phi=\{Q,[T_\mu,\phi]\}.
\]
这个等式意味着算符 $d\phi$ 是 $Q$-恰当的（$Q$-exact），
因而 $d\phi$ 在算符的 $Q$ 上同调类空间中的类为平凡。
这意味着：当我们改变 Riemann 曲面上的点 $x$ 时，局部算符 $\phi(x)$ 的 $Q$ 上同调类不变。
由此可知，在 $Q$ 上同调中这类 $Q$-闭局部算符彼此对易——只需把点移到类空间隔即可。
也就是说，算符的 $Q$ 上同调空间是一个 $\mathbb{Z}/2\mathbb{Z}$ 分次交换环。
我们总可以用 $0$-形式算符来代表这些算符。
这是「$Q$ 上同调类不依赖于度规」与「$Q$ 上同调类具有微分同胚不变性」这两件事的容易推论。
这些事实将在后面建立。

让我们回顾一个先前在第 11 讲 §11.3 中见过的构造。
给定一个 $Q$ 不变的 $0$-形式局部算符 $\phi=\phi^{(0)}$，我们定义一个 $1$-形式局部算符
\[
\phi^{(1)}_\mu=[T_\mu,\phi^{(0)}].
\]

% [Tutor Note: 作者的「简单计算」实为分次 Jacobi 恒等式：设 $Q$ 奇、$T_\mu$ 偶、$\phi$ 偶，则 $\{Q,[T_\mu,\phi]\}=[\{Q,T_\mu\},\phi]-[T_\mu,\{Q,\phi\}]$；只有当 $\{Q,\phi\}=0$（即 $\phi$ 为 $Q$-闭）时第二项才消失，从而得到 $d_\mu\phi=\{Q,[T_\mu,\phi]\}$。因此这条恒等式**只对 $Q$-闭算符**成立——作者紧接着就强调 $\phi$ 是 $Q$-closed，正是为此。若误把它当成对任意 $\phi$ 都成立的算符恒等式，就会错误地断言「所有」算符的上同调类都与位置无关。另外，要同时得到 $P_+$ 与 $P_-$ 的恰当性，(14.1) 中的 $Q$ 必须兼含 $Q_+$ 与 $Q_-$ 两种手征性；若只取单个 $Q_+$，则只有 $P_+$ 是恰当的，$\phi$ 便只能沿一个光锥方向平移，「上同调类与位置无关」也就不再成立。]

% ===== 原文 p.199 =====
以及
\[
\phi^{(1)}=\phi^{(1)}_\mu dx^\mu.
\]
与上一段所做的同样的计算表明
\[
d\phi^{(0)}=\{Q,\phi^{(1)}\}.
\]
这意味着：若 $\Sigma_1$ 是超 Riemann 曲面中的一个一维闭链（one-cycle），则
\[
\int_{\Sigma_1}\phi^{(1)}
\]
是 $Q$ 不变的。令 $\phi^{(2)}=[T,\phi^{(1)}]$。则
\[
d\phi^{(1)}=\{Q,\phi^{(2)}\}.
\]
因此 $\int_{\Sigma_2}\phi^{(2)}$ 是 $Q$ 不变的。
于是，以这种方式，从一个 $Q$ 不变的 $0$-形式局部算符 $\phi^{(0)}$ 出发，
我们可以构造出一族 $k$-形式局部算符 $\phi^{(k)}$，称为 $\phi^{(0)}$ 的后代，
从而得到更一般的 $Q$ 不变算符，其形式为
\[
\int_{\Sigma_k}\phi^{(k)},
\]
其中 $\Sigma_k$ 是 $\Sigma$ 中的一个 $k$ 维闭链。
我们的计算表明，对任何 $Q$ 不变的局部算符 $\phi^{(0)}$，
算符 $\phi^{(0)}+\phi^{(1)}+\phi^{(2)}$ 对总微分 $d-Q$ 是闭的。

这个「下降」（descent）程序在第 11 讲中也曾采用，那是在二维规范理论的背景下。

\subsection{扭曲理论以赋予其整体意义}
% 原文 §14.3：Twisting the theory to give it global meaning

到目前为止，我们一直在带固定坐标的 $\mathbb{R}^{2|4}$ 上工作。
如果我们旋转空间，则 $Q=Q_++Q_-$ 并非不变。
我们的目标是修改事物，直到 $Q$ 在旋转下不变，从而可以定义在更一般的欧氏 $2|4$ 维流形上。
让我们把欧氏旋转群记作 $SO(2)_E$。
欧氏理论中我们可支配的整个对称群是 $SO(2)_E\times U(1)_L\times U(1)_R$。
这个群到 $SO(2)_E$ 有显然的投影，并有一个回到第一个因子的分裂。
我们希望找到这个投影映射的另一个分裂，即乘积内部的一个群 $SO(2)'_E$，
它同构地投影到 $SO(2)_E$ 上，但具有与 $Q$ 对易这一额外性质。
容易看出，这样的分裂存在且唯一；它在 Lie 代数层面由
\[
K'=K\pm\frac12J_L\pm\frac12J_R
\]
给出，其中 $K$ 是 $SO(2)_E$ 的生成元，而符号取决于我们为 $Q$ 选定了四种可能性中的哪一种。
在我们集中关注的 $Q=Q_++Q_-$ 的情形，有
\[
K'=K+\frac12J_L-\frac12J_R.
\]
这些公式中出现 $1/2$ 这一事实意味着：扭曲实际上涉及 $SO(2)_E$ 的自旋形式（spin form）。

请注意，为了在 $n$ 维中做这样的扭曲，必须存在从 $\mathrm{Spin}(n)$ 到 $R$ 对称群的非平凡同态。
在二维，当 $R$ 对称群为 $U(1)$ 时（例如 $\mathcal{N}=1$ 超对称）就可以做到这一点。
但对 $n=4$，则必须有 $\mathcal{N}=2$ 超对称，

% [Tutor Note: 「容易看出这样的分裂存在且唯一」可以显式验证。设 $K'=K+aJ_L+bJ_R$，并要求 $[K',Q_++Q_-]=0$：由 $[K,Q_\pm]=\pm\frac12Q_\pm$、$[J_L,Q_-]=Q_-$、$[J_R,Q_+]=Q_+$（其余为零），$Q_+$ 分量给出 $\frac12+b=0$、$Q_-$ 分量给出 $-\frac12+a=0$，故 $a=\frac12$、$b=-\frac12$，即 $K'=K+\frac12J_L-\frac12J_R$，与文中一致；若改用 $Q_++\overline{Q}_-$，则 $b$ 变号而 $a$ 不变，得到 $K'=K+\frac12J_L+\frac12J_R$——这正是两种不等价扭曲（今称 A 模型与 B 模型）的来源。注意 $J_L,J_R$ 的荷归一化在这里必须取 $\pm1$ 而非 $\pm\frac12$，否则 $a,b$ 会差一个因子 2；这个归一化与 p.198 表格中「标准表示及其共轭」的说法相配。]

% [Tutor Note: 「涉及 $SO(2)_E$ 的自旋形式」这句话的后果常被低估：$K'$ 的半整数本征值意味着扭曲后的理论依赖自旋结构的选择，两种自旋结构给出的理论相差 $(-1)^F$——这也是为什么 p.200 要从「带选定自旋结构的 Riemann 曲面」出发。又，$n=4$ 时 $R$ 对称群必须是 $SU(2)\times U(1)$，因为只有 $SU(2)$ 才有来自 $\mathrm{Spin}(4)\cong SU(2)\times SU(2)$ 的非平凡同态；这一句是理解四维 Donaldson 理论如何作为 $\mathcal{N}=2$ 超对称规范理论的扭曲而出现的关键。]

% ===== 原文 p.200 =====
使得 $R$ 对称群为 $SU(2)\times U(1)$，而它才有来自 $\mathrm{Spin}(4)$ 的非平凡同态。

有了这些预备讨论，我们现在可以整体化我们关于局部算符的 $Q$ 上同调类的理论。
假设我们有一个局部以 $\mathbb{R}^{2|4}$ 为模型的 $2|4$ 维超 Riemann 曲面。
给定任何带有选定自旋结构的普通（玻色）Riemann 曲面，
都有相应的一个非常简单的超 Riemann 曲面：它是一个分裂的 $2|4$ 维复流形，其奇方向为
\[
\bigl(\Pi K^{1/2}\oplus\Pi\overline{K}^{1/2}\bigr)\oplus
\bigl(\Pi\overline{K}^{1/2}\oplus\Pi K^{1/2}\bigr).
\]
这里 $\theta^+,\overline{\theta}^+$ 是前两个线丛的截面，而 $\theta^-,\overline{\theta}^-$ 是后两个线丛的截面。
现在，我们将通过「扭曲」改变这个超流形的奇部。
我们把包含 $\theta^\pm$ 的自旋丛乘以典范丛（canonical bundle）平方根的 $\pm1/2$ 次幂来扭曲。
扭曲的结果是把上述线丛替换为
\[
\bigl(\Pi\mathcal{O}\oplus\Pi K\bigr)\oplus\bigl(\Pi\mathcal{O}\oplus\Pi\overline{K}\bigr),
\]
并使超对称
\[
Q=\frac{\partial}{\partial\theta^+}+\frac{\partial}{\partial\theta^-}
\]
成为一个整体对称性。
局部地，这个对称性与我们在 $\mathbb{R}^{2|4}$ 上一直讨论的对称性 $Q$ 一致。
任何一般地可以在超 Riemann 曲面上表述的 $\mathcal{N}=2$ 理论，
在被表述在我们刚刚描述的那类特殊超 Riemann 曲面上时，都将以 $Q$ 为对称性。
这保证了扭曲理论的整体版本存在；当然，它可以按分量场（component field）显式描述。

一旦我们把超对称 $Q$ 整体化，先前所有的讨论都可以整体化。
特别地，对一个 $Q$-闭局部算符 $\mathcal{O}(p)$，$\mathcal{O}(p)$ 的 $Q$ 上同调类与点 $p$ 无关。
于是，这个整体超 Riemann 曲面上的局部算符的 $Q$ 上同调类
生成一个 $\mathbb{Z}/2\mathbb{Z}$ 分次交换环，称为该理论的手征环。

\subsubsection{某些关联函数及其基本性质}
% 原文 §14.3.1：Certain Correlation Functions and Their Basic Properties

我们将考虑下列类型的路径积分：
\[
Z_{(\mathcal{O}_\alpha)}=\int D\Phi\,e^{-\mathcal{L}}
\prod_{\alpha=1}^{n}\int_{\Sigma_\alpha}\mathcal{O}_\alpha^{(k_\alpha)},
\]
其中对每个 $\alpha$，算符 $\mathcal{O}_\alpha$ 是 $Q$-闭局部算符，$k_\alpha=0,1,2$，
而 $\Sigma_\alpha$ 是超 Riemann 曲面 $\Sigma$ 中的一个 $k_\alpha$-维闭链。

以下是 $Z_{(\mathcal{O}_\alpha)}$ 的基本性质。

1). 它与超 Riemann 曲面上的度规无关。
理由是：当你把度规作 $\delta g$ 的变动时，你插入 $\int_\Sigma \delta g\cdot T$，其中 $T$ 是应力-能量张量。
但正如我们先前所见，$T$ 在如下意义上是 $Q$-恰当的：存在场 $S_{\mu\nu}$ 使 $T_{\mu\nu}=\{Q,S_{\mu\nu}\}$。
于是我们看到，这个插入产生一个形如 $\{Q,\cdot\}$ 的项。
这种项的路径积分为零，因为场是 $Q$ 不变的。
（还请注意到，必须假定 $Q$ 保持路径积分中的测度。若原理论是超对称的，
且流 $\frac12J_L-\frac12J_R$ 可被规范化，则情况确实如此。）

% [Tutor Note: 这里把「度规变动给出 $\int\delta g\cdot T$」当作显然，其实用到了两件事：(i) $T_{\mu\nu}$ 定义为作用量对度规的变分，即 $\delta S=\frac{1}{2}\int\sqrt{g}\,\delta g^{\mu\nu}T_{\mu\nu}$，与符号约定相差一个整体因子；(ii) 在二维，$T_{\mu\nu}$ 还有 Belinfante 改进（improvement）的任意性，而 $\{Q,S_{\mu\nu}\}$ 这一表述挑出了与超共形代数相容的那一个「改进后」的应力张量。另外，括号里「$Q$ 保持测度」这一点不可省：若该 $R$ 流有反常，路径积分测度在 $Q$ 下并不保持不变，$\{Q,\cdot\}$ 的积分就不再为零，本节所有结论都会失效——这正是前两节花大力气讨论 $k=0$ 的原因。]

% [Tutor Note: $(\Pi K^{1/2}\oplus\Pi\overline{K}^{1/2})\oplus(\Pi\overline{K}^{1/2}\oplus\Pi K^{1/2})$ 到 $(\Pi\mathcal{O}\oplus\Pi K)\oplus(\Pi\mathcal{O}\oplus\Pi\overline{K})$ 的替换，暗含了对两组奇坐标取**相反**的扭曲幂次：用 $\overline{K}=K^{-1}$ 可验，第一组 $(\theta^+,\overline{\theta}^+)$ 须乘 $K^{-1/2}$ 才能给出 $(\mathcal{O},K)$，而第二组 $(\theta^-,\overline{\theta}^-)$ 须乘 $K^{+1/2}$ 才能给出 $(\mathcal{O},\overline{K})$；若对两组用同一符号，就会得到 $K^{-1}$ 这种没有整体截面的丛，扭曲后的 $Q$ 也就不再整体定义。原文「to the power $\pm1/2$」中的 $\pm$ 正是此意，极易读漏。又，$\Pi$ 表示把纤维的奇偶性翻转（偶丛变奇丛），它保证 $\theta$ 是反对易的；$\mathcal{O}$ 指结构层。]

% ===== 原文 p.201 =====
（接上页：关联函数与度规无关这一事实的一个重要推论是：
$0$-形式算符 $\mathcal{O}^{(0)}_\alpha$ 必须是标量算符，即自旋为零；
自旋更高的对象在局部旋转下不会不变，它们在关联函数中的插入也就不可能与度规无关。）

2). 它与 $\mathcal{L}$ 中大多数参数无关。
首先，关联函数与任何形如 $\int d^2y\,d^4\theta(\cdots)$ 的项无关。
理由是：为计算这样的项，我们要对全部四个 $Q$ 求导，因此任何这样的项都具有 $\{Q,\dots\}$ 的形式。
带这种插入的积分因与上一段相同的理由而为零。
同理，取决于我们使用的是 $Q_+,Q_-,\overline{Q}_+,\overline{Q}_-$ 中的哪个组合，
项
\[
\int d^2y\,d\theta^+\,d\theta^-\,W+{\rm c.c.}
\]
与
\[
\int d^2y\,d\overline{\theta}^+\,d\theta^-\,\widetilde{W}+{\rm c.c.}
\]
中只有一个可能是非平凡的。
事实上，关联函数与其中之一无关，而对另一个全纯地变化。

3). 我们可以用施于 $Q$ 的作用的不动点定理来计算 $Z_{(\mathcal{O}_\alpha)}$。
我们把 $Q$ 看作（无穷小地）作用在理论中所有场的函数空间上。
若这个作用没有不动点，则路径积分的值为零。
一般地，我们可以把计算归结为对 $Q$ 的不动点集的积分。
这些当然就是那些在超对称 $Q$ 下不变的场位形。

\subsubsection{态的 $Q$ 上同调}
% 原文 §14.3.2：Q-cohomology of States

现在我们从局部算符、路径积分与关联函数的世界，转到 Hamilton 框架。
我们的 Riemann 曲面现在是柱面 $S^1\times\mathbb{R}$。
我们把所有相关的自旋丛在 $S^1$ 上平凡化。
我们有圆周上量子理论的态空间 $\mathcal{H}$，以及 $\mathcal{H}$ 上与时间平移相联系的 Hamilton 算符 $H$。
由于 $Q$ 作用于理论，它在 $\mathcal{H}$ 上产生一个算符（我们也称之为 $Q$），与 $H$ 对易。
此外，由超代数中的对易关系我们看到
\[
H=\{Q,\widetilde{Q}\},
\]
其中 $\widetilde{Q}$ 是 $Q$ 的伴随。当然，我们仍有关系 $Q^2=0$，
因此可以取 $Q$ 对 $\mathcal{H}$ 的作用的上同调：
\[
H^*(Q)=\ker Q/\operatorname{Im}Q.
\]
若我们的理论足够好，使得 $H$ 具有离散谱，则 $H^*(Q)$ 可等同于 $\{\psi\,|\,H\psi=0\}$。
（我们稍后要讨论的规范理论例子并非如此，但对 $\sigma$ 模型成立。）
一般地，$H^*(Q)$ 按 $U(1)\times U(1)$ $R$ 对称性群的作用而双分次。

% [Tutor Note: 两个 F 项的测度并不相同：第一个是通常的 $\int d^2y\,d\theta^+d\theta^-W$，第二个是「扭曲」F 项 $\int d^2y\,d\overline{\theta}^+d\theta^-\widetilde{W}$（注意是 $\overline{\theta}^+$ 配 $\theta^-$，且超势带波浪号），它正是第 13 讲中扭曲超势 $\widetilde{\mathcal{W}}_{\rm eff}$ 所对应的项。千万不要因为「看起来不对称」而擅自改成 $d\overline{\theta}^+d\overline{\theta}^-$，否则该项将不再是对任何 $Q$ 都恰当的恰当形式。一般地，$\int d^4\theta$ 型的项对四个超荷全部求导，故对四种 $Q$ 都是 $\{Q,\cdot\}$；而两个 F 项各自只对其中两个超荷恰当，所以对给定的 $Q$ 只有一个能存活，另一个则全纯地依赖于参数——这正是作者说「一个无关、另一个全纯变化」的含义。]

% [Tutor Note: $H=\{Q,\widetilde{Q}\}$ 中的 $\widetilde{Q}$ 是 $Q$ 的伴随，但它**不是** $Q$ 本身，也未必是同一扭曲理论所保持的超荷：若 $Q=Q_++Q_-$，则 $Q^\dagger=\overline{Q}_++\overline{Q}_-$，这正是 §14.2 列举的第三种组合，而两种扭曲理论一般并不等价。扭曲理论通常还是非正定的（$K'$ 不是任何算符的伴随），故 $H$ 未必有下界；把 $H^*(Q)$ 等同于零能态除了 $Q^2=0$ 之外，还需要「$H$ 有离散谱」这一假设（即 Hodge 型论证中的椭圆性与紧性）。作者随即声明规范理论的例子不满足此假设，原因正在于平坦联络模空间非紧、$H$ 的谱连续。]
